% Deployment walkthrough for the nosnitch package.
% Requires: \usepackage{listings}, \usepackage{xcolor} (or similar) in the
% including document. The listings below use a generic shell style; adjust
% \lstset to match your paper's conventions.

\section{Deployment}
\label{sec:deployment}

The \texttt{nosnitch} package implements the AirSnitch mitigations across
two layers: a userspace guard daemon (shipped as an OpenWrt package) and
a set of source-level patches to \texttt{hostapd}. The daemon covers
Patches~2 and~3 in full. The hostapd patches are required for Patch~1's
core primitives (broadcast filtering, per-client GTK in secondary
handshakes) and for Patch~3's cross-BSSID duplicate-MAC rejection. Two
deployment paths are supported:

\begin{description}
  \item[Full deployment (\S\ref{sec:deploy-full}).] Rebuild OpenWrt from
  source with the hostapd patches applied. All mitigations active.
  \item[Runtime-only deployment (\S\ref{sec:deploy-runtime}).] Install
  the package on a stock OpenWrt router via \texttt{opkg}. Patches~2
  and~3 (userspace) are active; Patch~1 primitives requiring hostapd
  source changes are not.
\end{description}

Full deployment is required to reproduce the results reported in this
paper. Runtime-only deployment is documented as a defense-in-depth
fallback for operators who cannot maintain a custom firmware image.

\subsection{Full Deployment}
\label{sec:deploy-full}

\subsubsection{Prerequisites}

\begin{itemize}
  \item A Linux build host with the OpenWrt toolchain prerequisites
  (see the OpenWrt Developer Guide).
  \item The target device's OpenWrt version and profile. Reference
  deployment: GL.iNet GL-MT6000 (Flint 2), OpenWrt 24.10,
  profile \texttt{mediatek/filogic}.
  \item Approximately 20\,GB of free disk for the buildroot.
\end{itemize}

\subsubsection{Step 1: Obtain the OpenWrt buildroot}

\begin{lstlisting}[language=bash]
git clone https://git.openwrt.org/openwrt/openwrt.git
cd openwrt
git checkout v24.10.0
./scripts/feeds update -a
./scripts/feeds install -a
\end{lstlisting}

\subsubsection{Step 2: Register the nosnitch feed}

Append the following line to \texttt{feeds.conf.default} (or create
\texttt{feeds.conf}):

\begin{lstlisting}
src-git nosnitch https://github.com/JamesLove03/NoSnitch.git
\end{lstlisting}

\noindent Install the feed:

\begin{lstlisting}[language=bash]
./scripts/feeds update nosnitch
./scripts/feeds install nosnitch
\end{lstlisting}

\subsubsection{Step 3: Apply the hostapd patches}
\label{sec:deploy-hostapd-patches}

The nosnitch feed ships a helper script that installs the hostapd
patches into the buildroot in one invocation. It distinguishes between
the two patch flavours automatically: source-level patches against
upstream hostapd are copied into
\texttt{package/network/services/hostapd/patches/} for OpenWrt's build
system to apply during compile, while the buildroot-level patch against
\texttt{hostapd.sh} is applied in-tree with \texttt{patch~-p1
--forward}. The \texttt{--forward} flag makes the script idempotent;
re-running on an already-patched tree is a no-op.

A POSIX-shell version and a Python version are provided; both behave
identically. Pick whichever is more convenient on the build host:

\begin{lstlisting}[language=bash]
feeds/nosnitch/nosnitch/scripts/install-hostapd-patches.sh .
# or:
feeds/nosnitch/nosnitch/scripts/install-hostapd-patches.py .
\end{lstlisting}

\noindent The single positional argument is the OpenWrt source root
(here, the current directory). The script validates that the target
looks like an OpenWrt tree before touching anything.

The source-level patches target hostapd~2.11. On other versions the
hunks may fail with fuzz; resolve any \texttt{.rej} files by
re-anchoring to the current line numbers before proceeding to Step~5.

\subsubsection{Step 4: Select packages and target}

\begin{lstlisting}[language=bash]
make menuconfig
\end{lstlisting}

\noindent In the menu:

\begin{enumerate}
  \item Set \emph{Target System} and \emph{Target Profile} to the
  deployment device.
  \item Under \emph{Network}, mark \texttt{nosnitch} as \texttt{<*>}
  (built-in) or \texttt{<M>} (loadable module).
  \item Exit and save.
\end{enumerate}

\subsubsection{Step 5: Build}

\begin{lstlisting}[language=bash]
make -j$(nproc) package/network/services/hostapd/{clean,compile} V=s
make -j$(nproc) package/feeds/nosnitch/nosnitch/{clean,compile} V=s
make -j$(nproc)
\end{lstlisting}

\noindent The resulting firmware image is placed under
\texttt{bin/targets/<target>/<profile>/}.

\subsubsection{Step 6: Flash the device}

For a device with OpenWrt already running, use \texttt{sysupgrade}:

\begin{lstlisting}[language=bash]
scp bin/targets/mediatek/filogic/openwrt-*-sysupgrade.bin \
    root@192.168.1.1:/tmp/
ssh root@192.168.1.1 "sysupgrade -v /tmp/openwrt-*-sysupgrade.bin"
\end{lstlisting}

\noindent For a stock-firmware device, follow the vendor's recovery
procedure with the corresponding \texttt{-factory} image.

\subsubsection{Step 7: Configure wireless}

After reboot, SSH into the router and add the following options to each
\texttt{config wifi-iface} block in \texttt{/etc/config/wireless} that
the mitigations should protect:

\begin{lstlisting}
option isolate              '1'
option multicast_to_unicast '1'
option proxy_arp            '1'
option wpa_strict_rekey     '1'
option wpa_group_rekey      '600'
option disable_dgaf         '1'
option per_client_gtk       '1'
\end{lstlisting}

\noindent Patch~3 Step~5 (\texttt{ieee80211w=2}, MFP Required) is
defaulted by the \texttt{950-hostapd-sh-uci.patch} when no
\texttt{ieee80211w} is set on the wifi-iface; no manual setting is
required. To keep legacy clients (Win7, very old Android, some IoT)
working, explicitly set \texttt{option ieee80211w '0'} or \texttt{'1'}
on the affected interface.

For IPv6-serving networks, add \texttt{option ndp 'proxy'} to the
corresponding \texttt{config interface} block in
\texttt{/etc/config/network}.

\subsubsection{Step 8: Configure the guard daemon}

The daemon reads \texttt{/etc/config/nosnitch}. The shipped defaults are
correct for a single-AP deployment on \texttt{br-lan}. Verify or adjust:

\begin{lstlisting}
config nosnitch 'main'
    option enabled           '1'
    option l3_isolation      '1'
    option mac_lock          '1'
    option mac_lock_bridge   '1'
    option mac_lock_network  '0'
    option bridge            'br-lan'
    option sync_interface    'eth0'
\end{lstlisting}

\noindent Set \texttt{mac\_lock\_network} to \texttt{'1'} only in
multi-AP deployments where \emph{every} AP on the wired LAN runs
\texttt{nosnitch}. Mixed deployments will generate false-positive
spoofing events when a legitimate client roams to an unpatched AP.

\subsubsection{Step 9: Apply configuration}

\begin{lstlisting}[language=bash]
wifi reload
/etc/init.d/network restart
fw4 reload
sysctl -p /etc/sysctl.d/30-nosnitch.conf
/etc/init.d/nosnitch enable
/etc/init.d/nosnitch restart
\end{lstlisting}

\subsubsection{Step 10: Verify}

\begin{lstlisting}[language=bash]
logread -e nosnitch
nft list set inet fw4 wifi_clients
ebtables -L FORWARD
hostapd_cli -i wlan0 status | grep -E \
    'per_client_gtk|proxy_arp|multicast_to_unicast'
\end{lstlisting}

\noindent A successful deployment shows a \texttt{started} line in the
log output, an (initially empty) nftables set, no lingering ebtables
rules, and the hostapd options reported as enabled. End-to-end
validation uses the AirSnitch test tool; each attack primitive must
fail when run from a client on the protected SSID.

\subsection{Runtime-Only Deployment}
\label{sec:deploy-runtime}

On a router already running stock OpenWrt, the \texttt{nosnitch}
package can be installed without rebuilding firmware. Patch~1's
hostapd-resident primitives will not be active; the remaining
mitigations are.

\subsubsection{Step 1: Build the package}

Using the OpenWrt SDK for the target platform:

\begin{lstlisting}[language=bash]
cd openwrt-sdk-*
./scripts/feeds update nosnitch
./scripts/feeds install nosnitch
make package/nosnitch/compile V=s
\end{lstlisting}

\noindent The \texttt{.ipk} is produced under \texttt{bin/packages/}.
The hostapd-patches helper script is not invoked in this path; the
stock hostapd binary is used as-is.

\subsubsection{Step 2: Install}

\begin{lstlisting}[language=bash]
scp bin/packages/*/nosnitch/nosnitch_*.ipk root@router:/tmp/
ssh root@router "opkg update && opkg install /tmp/nosnitch_*.ipk"
\end{lstlisting}

\subsubsection{Step 3: Configure and reload}

Apply the wireless and nosnitch configuration from
\S\ref{sec:deploy-full} Steps~7--9, omitting the
\texttt{per\_client\_gtk} option (the unpatched hostapd does not
recognize it). The MFP default from the \texttt{950} patch is also not
in effect; if MFP is desired, set \texttt{option ieee80211w '2'}
explicitly on each protected wifi-iface.

\subsection{Operational Notes}

\begin{itemize}
  \item \textbf{Patch drift.} The hostapd patches carry approximate
  context. After every OpenWrt point release, rebase the patch files
  against the new hostapd source and rerun the helper script
  (\S\ref{sec:deploy-hostapd-patches}). Because the script uses
  \texttt{patch~--forward}, it tolerates a partially-applied tree from
  a previous build.
  \item \textbf{Per-radio interface enumeration.} The daemon's
  cross-radio MAC locking requires the wireless interface list to be
  populated. On devices with non-default interface names, set
  \texttt{option wlan\_ifaces 'wlan0 wlan1'} in
  \texttt{/etc/config/nosnitch}.
  \item \textbf{Gateway MAC locking.} Patch~3 Step~3 (preventing
  gateway MAC spoofing) requires the upstream gateway's MAC address.
  Set \texttt{option gateway\_mac} in \texttt{/etc/config/nosnitch} for
  deployments behind a known upstream router.
  \item \textbf{Uninstall.} \texttt{opkg remove nosnitch} removes the
  daemon, the nftables ruleset fragment, and the sysctl overlay. The
  hostapd patches remain baked into the firmware image; reverting
  them requires a rebuild.
\end{itemize}
